% abstract.tex -- input by main.tex
\section*{Abstract}
Virus Isolation in cell culture commonly pairs two manipulations: inoculation and a step-down in fetal bovine serum (FBS) from a growth medium (typically 10\%) to a maintenance medium (typically 2\%), followed by reading cytopathic effect (CPE) as evidence of virus. We asked whether the serum step-down alone is sufficient to produce morphology that trained observers describe as CPE. (i) In a dual-FBS corpus of 170 Isolation papers, the modal recipe was 10\% growth to 2\% maintenance (medians 10.0\% and 2.0\%; plotted $n=169$), matching eight institutional guidelines. (ii) Isolation-stage negative controls with matched maintenance medium were reported in 11 of 173 scored papers (Tier A); 110 of 173 reported no qualifying control (Tier D). (iii) In a contract-laboratory experiment with uninoculated Vero E6 cells held at 10\% versus 2\% FBS, the laboratory's own image descriptions named CPE-type features on 1 of 9 described 10\%-FBS frames and 10 of 13 described 2\%-FBS frames; two independent raters blind to culture condition each scored more CPE-type morphology on the 2\%-FBS culture (mean CPE-item score 0.87 vs 1.65 and 2.41 vs 4.91). \assumed{Records obtained under a public-records request indicate that reference-laboratory protocols ... [replace with the FOIA result].} These results show that a medium change routinely applied at inoculation can by itself produce CPE-classified morphology, and that published Isolation methods rarely include the control that would separate the two. \todoA{Final sentence on implications once FOIA and IR3v Pass 2 are in.}
